\section*{الفصل \ref{ch:b1:central-forces} --- القوى المركزية: الكواكب والأقمار الصناعية}

\begin{solution}{exo:b1:central-forces:1}
$r = \qty{6770}{km}$: ومنه $v = \sqrt{GM/r} = \qty{7.7}{km/s}$؛ و $T = 2\pi r/v
= \qty{5.5e3}{s} = \qty{92}{min}$؛ و $1440/92 \approx 16$ شروقًا في اليوم.
\end{solution}

\begin{solution}{exo:b1:central-forces:2}
$r = (GMT^2/4\pi^2)^{1/3} = (3.99\times10^{14} \times 7.42\times10^9/39.5)^{1/3}
= \qty{4.22e7}{m}$: أي ارتفاع \qty{35800}{km}؛ و $v = 2\pi r/T = \qty{3.1}{km/s}$.
\end{solution}

\begin{solution}{exo:b1:central-forces:3}
$v_e = \sqrt{2GM/R}$: فللأرض \qty{11.2}{km/s}؛ وللقمر $\sqrt{2 \times 4.90\times
10^{12}/1.74\times10^6} = \qty{2.4}{km/s}$. ولجزيئات الغاز عند بضع مئات
الأمتار في الثانية ذيل في توزيع سرعاتها فوق
\qty{2.4}{km/s}؛ فعلى مدى الأزمنة الجيولوجية تسرّب غاز القمر.
\end{solution}

\begin{solution}{exo:b1:central-forces:4}
$T = a^{3/2}$ (بالسنوات والوحدات الفلكية): فللمريخ \qty{1.88}{yr}؛ وللمشتري \qty{11.9}{yr}؛
وعند \qty{30}{AU}: \qty{164}{yr}.
\end{solution}

\begin{solution}{exo:b1:central-forces:5}
$\Delta E = \tfrac12 mv^2 + GMm(1/R - 1/r) = 2.94\times10^{10} + 3.99\times
10^{17} \times (1.570 - 1.477)\times10^{-7} = 2.94\times10^{10} + 0.37\times
10^{10} = \qty{3.3e10}{J}$ --- أي \qty{33}{MJ/kg}، وهي طاقة \qty{0.7}{kg}
من البنزين لكل كيلوغرام يوضع في المدار (وتحتاج الصواريخ إلى أكثر من ذلك بكثير،
لأن عليها رفع وقودها).
\end{solution}

\begin{solution}{exo:b1:central-forces:6}
$a = (0.586 + 35.1)/2 = \qty{17.8}{AU}$؛ و $e = (35.1 - 0.586)/35.7 = 0.967$؛
و $T = 17.8^{3/2} = \qty{75}{yr}$. و $v_p^2 = GM(2/r_p - 1/a)$ مع $r_p =
8.77\times10^{10}$ م و $a = 2.67\times10^{12}$ م: ومنه $v_p = \qty{54}{km/s}$؛
و $v_a = v_pr_p/r_a = \qty{0.91}{km/s}$.
\end{solution}

\begin{solution}{exo:b1:central-forces:7}
$a = \qty{24490}{km}$. و $v_1 = \qty{7.67}{km/s}$؛ وعلى الإهليلج عند $r_1$:
$\sqrt{GM(2/r_1 - 1/a)} = \qty{10.07}{km/s}$: ومنه $\Delta v_1 = \qty{2.4}{km/s}$.
وعند $r_2$: $\sqrt{GM(2/r_2 - 1/a)} = \qty{1.62}{km/s}$، والسرعة الدائرية
\qty{3.07}{km/s}: ومنه $\Delta v_2 = \qty{1.45}{km/s}$. والمجموع \qty{3.85}{km/s}؛
والمدة $\pi\sqrt{a^3/GM} = \qty{1.9e4}{s} = \qty{5.3}{h}$.
\end{solution}

\begin{solution}{exo:b1:central-forces:8}
$E_{p,\mathrm{eff}}' = GMm/r^2 - L^2/mr^3 = 0$ عند $r_c = L^2/GMm^2$.
و $E_{p,\mathrm{eff}}'' = -2GMm/r^3 + 3L^2/mr^4$؛ وعند $r_c$ تعطي $L^2 = GMm^2r_c$
القيمةَ $E'' = GMm/r_c^3$، ومنه $\omega_r^2 = E''/m = GM/r_c^3 = \omega_{\text{مداري}}^2$.
فيهتزّ اضطراب قطري صغير مرة واحدة بالضبط في كل دورة: أي إن
المدار ينغلق على نفسه --- فهو إهليلج قليل \omterm{thm:b1:central-forces:conics}{الاختلاف المركزي}، لا وردة.
\end{solution}

\begin{solution}{exo:b1:central-forces:9}
لا تؤثر إلا الثقالة في المحطة والطاقم معًا؛ وكلاهما يسقط على
المدار نفسه، فلا حاجة إلى قوة تماسّ بينهما --- فلا أرضية تدفع،
ومن هنا ``انعدام الوزن''. و $g = GM/r^2 = 3.99\times10^{14}/(6.77\times10^6)^2
= \qty{8.7}{m/s^2}$: فالثقالة بكامل قوتها تقريبًا؛ والوزن ليس
غائبًا، بل هو منفَق كله في ثني المسار.
\end{solution}

\begin{solution}{exo:b1:central-forces:10}
$a = 4\pi^2r/T^2 = 39.5 \times 3.84\times10^8/(2.36\times10^6)^2 =
\qty{2.7e-3}{m/s^2}$؛ و $g(R/r)^2 = 9.81/60.3^2 = \qty{2.7e-3}{m/s^2}$: أي إن
قانون القوة نفسه الذي يُسقط تفاحة يمسك القمر، مخفَّفًا بالمقدار $1/r^2$.
\end{solution}

\begin{solution}{exo:b1:central-forces:11}
$E_k = 2Ze^2/4\pi\varepsilon_0r_{\min}$: ومنه $r_{\min} = 8.99\times10^9 \times 2 \times 79
\times (1.60\times10^{-19})^2/8.0\times10^{-13} = \qty{4.5e-14}{m} = \qty{45}{fm}$،
أي ستة أضعاف نصف القطر النووي: فيرتدّ الجسيم $\alpha$ قبل الملامسة.
ولبلوغ \qty{7}{fm}: $E_k \approx \qty{32}{MeV}$.
\end{solution}

\begin{solution}{exo:b1:central-forces:12}
$E_m = -GMm/2r$ يتناقص $\Rightarrow$ يتناقص $r$ $\Rightarrow$ فتكبر $v =
\sqrt{GM/r}$. فترتفع \omterm{thm:b1:work-and-energy:ke}{الطاقة الحركية} $GMm/2r$ بالمقدار $\abs{\Delta E}$
بينما تهبط \omterm{def:b1:work-and-energy:conservative}{الطاقة الكامنة} $-GMm/r$ بالمقدار $2\abs{\Delta E}$: أي إن نصف
\omterm{def:b1:work-and-energy:conservative}{الطاقة الكامنة} المحرَّرة يصير حركية، والنصف الآخر يتبدد حرارةً
بفعل السحب.
\end{solution}

\begin{solution}{pb:b1:central-forces:1}
\textbf{1.} $T = \qty{43082}{s}$: ومنه $r = (GMT^2/4\pi^2)^{1/3} = (3.986\times10^{14}
\times 1.856\times10^9/39.48)^{1/3} = \qty{2.656e7}{m}$؛ والارتفاع
\qty{20190}{km}.

\textbf{2.} $v = \sqrt{GM/r} = \qty{3.87}{km/s}$.

\textbf{3.} $a = v^2/r = GM/r^2 = \qty{0.565}{m/s^2} = g(R/r)^2$، أي أضعف سبع عشرة
مرة مما هو عند السطح --- لكنه ليس معدومًا: فالقمر الصناعي يسقط
باستمرار حول الأرض.

\textbf{4.} $E_m/m = -GM/2r = \qty{-7.5}{MJ/kg}$؛ و $E_k/m = \qty{+7.5}{MJ/kg}$ و
$E_p/m = \qty{-15.0}{MJ/kg}$.

\textbf{5.} تدور الأرض دورة واحدة في كل يوم نجمي بالنسبة إلى النجوم
(وهي مرجع المدار)؛ وبعد دورتين يعود القمر الصناعي فوق
النقطة الأرضية نفسها: فتتكرر المسارات الأرضية كل يوم نجمي.

\textbf{6.} $\sin\beta = R/r = 0.240$: ومنه $\beta = 13.9^\circ$، أي قرص
عرضه $28^\circ$.

\textbf{7.} \omterm{prop:b1:central-forces:escape}{المدار الثابت بالنسبة إلى الأرض} \qty{42200}{km}: أي أعلى بالعامل $1.59$؛
فمدار الملاحة دونه.

\textbf{8.} $GM(1/R - 1/2r) = 3.986\times10^{14}(1.570\times10^{-7} - 1.88\times
10^{-8}) = \qty{5.5e7}{J/kg} = \qty{55}{MJ/kg}$.

\textbf{9.} $r_1 = \qty{6771}{km}$: $v = \qty{7.67}{km/s}$.

\textbf{10.} $a = \qty{16665}{km}$؛ وعند $r_1$: $\sqrt{GM(2/r_1 - 1/a)} =
\qty{9.69}{km/s}$، ومنه $\Delta v_1 = \qty{2.0}{km/s}$؛ وعند $r_2$: \qty{2.47}{km/s}
مقابل الدائرية \qty{3.87}{km/s}، ومنه $\Delta v_2 = \qty{1.4}{km/s}$؛ والمجموع
\qty{3.4}{km/s}.

\textbf{11.} $\pi\sqrt{a^3/GM} = \qty{1.07e4}{s} \approx \qty{3.0}{h}$.

\textbf{12.} الإفلات من المدار المنخفض: $(\sqrt2 - 1)\times 7.67 = \qty{3.2}{km/s}$:
فقيمة المهمة \qty{3.4}{km/s} تتجاوزها --- أي إن بلوغ هذا المدار يكلّف
نحو ما تكلّفه مغادرة الأرض كلها.

\textbf{13.} $m_0/m_1 = \eu^{3.4/3.0} = 3.1$: ومنه $m_1 = \qty{640}{kg}$،
والوقود الدافع \qty{1360}{kg} --- أي ثلثا الكتلة.

\textbf{14.} $\cos\theta = 0.240$ و $\theta = 76^\circ$؛ والنسبة $(1 - 0.240)/2
= 0.38$.

\textbf{15.} $24 \times 0.38 = 9$ في المتوسط؛ والأربعة متجاوَزة
بأريحية في كل وقت تقريبًا.

\textbf{16.} في السمت: $2.02\times10^7/c = \qty{67}{ms}$؛ وعند الأفق: البعد
$\sqrt{r^2 - R^2} = \qty{2.58e7}{m}$، أي \qty{86}{ms}.

\textbf{17.} $3/c = \qty{10}{ns}$.

\textbf{18.} انزياح ساعة المستقبِل مجهول رابع إلى جانب
إحداثياته الثلاثة: فأربع مسافات، وأربع معادلات.

\textbf{19.} $10^{-8}\ \unit{s}/\qty{86400}{s} \approx \num{1e-13}$.

\textbf{20.} $v^2/2c^2 = (3874)^2/(2 \times 8.99\times10^{16}) = \num{8.3e-11}$:
ومنه $8.3\times10^{-11} \times 86400 = \qty{7.2}{\micro s}$ في اليوم، تباطؤًا.

\textbf{21.} $GM(1/R - 1/r)/c^2 = 3.986\times10^{14} \times 1.19\times10^{-7}/
8.99\times10^{16} = \num{5.3e-10}$: أي \qty{45.7}{\micro s} في اليوم، تسارعًا.

\textbf{22.} المحصلة $+\qty{38.5}{\micro s}$ في اليوم؛ و $c \times 38.5\times10^{-6}
= \qty{11.5}{km}$ من الخطأ في اليوم.

\textbf{23.} $(10.23 - 10.22999999543)/10.23 = \num{4.47e-10}$، وهي توافق
المحصلة $4.5\times10^{-10}$ (أي \qty{38.5}{\micro s} على \qty{86400}{s}).

\textbf{24.} انتفاخ الأرض الاستوائي، وجذب القمر والشمس،
وضغط الإشعاع الشمسي (وإشعالات صغيرة للدوافع): فيقيس
القطاع الأرضي المدارات باستمرار ويرفع تقاويم فلكية جديدة
تبثّها الأقمار.

\textbf{25.} تحدّد الطاقة \omterm{def:b1:angular-momentum:L}{والعزم الحركي} مدارًا دائريًّا انطلاقًا من
دوره؛ ويضعه كبلر الثالث؛ وتحدّد الهندسة التغطية وأزمنة الضوء؛
أما الساعات --- وهي قلب المنظومة --- فيجب تصحيحها من أجل
النسبية بمقدار \qty{38}{\micro s} في اليوم، وإلا انحرف التحديد بعشرة
كيلومترات مع حلول الليل.
\end{solution}
